% !TEX root = ../../../main.tex
% 所属：第二部分：多变量分析学 / 第2章：多元函数微分学
% 节：2.1 偏导数,全微分
% 正文按原文件顺序保存；由 ../chapter2.tex 调用。
\section{偏导数,全微分}

\subsection*{偏导数}

以二元函数
\[
z=f(x,y)
\]
为例.

\begin{definition}
	对 $x$ 的偏导数定义为
\[
\frac{\partial z}{\partial x}\left(\text{或 }\frac{\partial f}{\partial x}\right)
:=\lim_{\Delta x\to0}\frac{f(x+\Delta x,y)-f(x,y)}{\Delta x}.
\]
类似地，对 $y$ 的偏导数定义为
\[
\frac{\partial z}{\partial y}\left(\text{或 }\frac{\partial f}{\partial y}\right)
:=\lim_{\Delta y\to0}\frac{f(x,y+\Delta y)-f(x,y)}{\Delta y}.
\]
若 $z=f(x,y)$ 在一点处对 $x,y$ 的偏导数 $\dfrac{\partial f}{\partial x}$、$\dfrac{\partial f}{\partial y}$ 都存在，则称 $z=f(x,y)$ 在该点\textbf{可偏导}.

\end{definition}
偏导数存在并不意味着函数在该点连续.

\begin{example}
	令
\[
f(x,y)=
\begin{cases}
0,&(x,y)=(0,0),\\[4pt]
\dfrac{xy}{x^2+y^2},&(x,y)\ne(0,0).
\end{cases}
\]
在原点处，有
\[
\frac{\partial f}{\partial x}(0,0)=f_x(0,0)
=\lim_{\Delta x\to0}\frac{f(\Delta x,0)-f(0,0)}{\Delta x}=0,
\qquad
\frac{\partial f}{\partial y}(0,0)=f_y(0,0)=0.
\]
$\dfrac{\partial f}{\partial x}$、$\dfrac{\partial f}{\partial y}$ 在原点处存在，但是$f$ 在原点处不连续.
\end{example}
\subsection*{方向导数}

仍以二元函数 $z=f(x,y)$ 为例.

\begin{definition}
	沿平面 $\R^2$ 中某方向
\[
\vect v=(\cos\alpha,\sin\alpha)
\]
的方向导数定义为
\[
\frac{\partial f}{\partial\vect v}(x_0,y_0)
=\lim_{t\to0}
\frac{f(x_0+t\cos\alpha,y_0+t\sin\alpha)-f(x_0,y_0)}{t}.
\]
若上式极限存在，则称 $f$ 在 $(x_0,y_0)$ 处沿 $\vect v$ 方向的方向导数存在，并记为 $\dfrac{\partial f}{\partial\vect v}$.
\end{definition}

若 $\dfrac{\partial f}{\partial x}$ 存在，取 $\vect v=(1,0)$，则
\[
\frac{\partial f}{\partial\vect v}
=\lim_{t\to0}\frac{f(x_0+t,y_0)-f(x_0,y_0)}{t}
=\left.\frac{\partial f}{\partial x}\right|_{(x_0,y_0)}.
\]
而当 $\vect v=(-1,0)$ 时，
\[
\frac{\partial f}{\partial\vect v}=-\frac{\partial f}{\partial x}.
\]


	\textbf{注：}偏导数存在时，方向导数不一定存在.同样地，沿任意方向的方向导数均存在，也不一定推出可微.


同样，若将 $\R^n$ 中的单位向量 $\vect v$（即满足 $|\vect v|=1$ 的向量）视为一个方向，就可类似定义 $n$ 元函数的方向导数.

\begin{definition}
	设 $D\subset\R^n$ 为开集，$\vect x^{\,0}=(x_1^0,x_2^0,\cdots,x_n^0)$ 为 $D$ 中一点，$\vect v=(v_1,v_2,\cdots,v_n)$ 为一方向.定义 $D$ 上的 $n$ 元函数
	\[
	u=f(x_1,x_2,\cdots,x_n)
	\]
	在点 $\vect x^{\,0}$ 沿方向 $\vect v$ 的方向导数为
	\[
	\frac{\partial f}{\partial\vect v}(x_1^0,x_2^0,\cdots,x_n^0)
	=\lim_{t\to0}
	\frac{f(x_1^0+tv_1,x_2^0+tv_2,\cdots,x_n^0+tv_n)-f(x_1^0,x_2^0,\cdots,x_n^0)}{t},
	\]
	如果等式右边的极限存在.
\end{definition}

\subsection*{全微分}

\begin{definition}
	对 $z=f(x,y)$，考虑全增量
\[
\Delta z=f(x+\Delta x,y+\Delta y)-f(x,y).
\]
若存在与 $\Delta x,\Delta y$ 无关的函数 $A(x,y)$、$B(x,y)$，使得
\[
\Delta z=A(x,y)\Delta x+B(x,y)\Delta y
+o\!\left(\sqrt{\Delta x^2+\Delta y^2}\right),
\]
则称 $f(x,y)$ 在该点\textbf{可微}，并称
\[
\dd z=A(x,y)\dd x+B(x,y)\dd y
\]
为 $f(x,y)$ 的\textbf{全微分}.
\end{definition}

如果 $f(x,y)$ 是可微的，则容易推出 $f(x,y)$ 是连续的.偏导数存在不一定推出连续，因此也不一定推出可微.

\begin{proposition}
若 $f(x,y)$ 在一点处可微，则 $f(x,y)$ 在该点可偏导，且
\[
\dd z=\frac{\partial f}{\partial x}\dd x+\frac{\partial f}{\partial y}\dd y.
\]
\end{proposition}

\begin{proof}
由可微的定义，
\[
f(x+\Delta x,y+\Delta y)-f(x,y)
=A(x,y)\Delta x+B(x,y)\Delta y
+o\!\left(\sqrt{\Delta x^2+\Delta y^2}\right).
\]
令 $\Delta y=0$，则
\[
\frac{f(x+\Delta x,y)-f(x,y)}{\Delta x}
=A(x,y)+\frac{o(|\Delta x|)}{\Delta x}.
\]
当 $\Delta x\to0$ 时，右端趋于 $A(x,y)$，故
\[
A(x,y)=\frac{\partial f}{\partial x}.
\]
同理，$B(x,y)=\dfrac{\partial f}{\partial y}$，从而结论成立.
\end{proof}

\begin{proposition}
	若 $f$ 在 $(x_0,y_0)$ 处可微，则 $f$ 在该点沿任意方向
	\[
	\vect v=(\cos\alpha,\sin\alpha)
	\]
	的方向导数均存在，且
	\[
	\frac{\partial f}{\partial\vect v}(x_0,y_0)
	=
	\frac{\partial f}{\partial x}(x_0,y_0)\cos\alpha
	+
	\frac{\partial f}{\partial y}(x_0,y_0)\sin\alpha.
	\]
\end{proposition}

\begin{proof}
	由 $f$ 在 $(x_0,y_0)$ 处可微，当 $t\to0$ 时，
	\[
	\begin{aligned}
	&f(x_0+t\cos\alpha,y_0+t\sin\alpha)-f(x_0,y_0)\\
	&\qquad=
	\frac{\partial f}{\partial x}(x_0,y_0)t\cos\alpha
	+
	\frac{\partial f}{\partial y}(x_0,y_0)t\sin\alpha
o(t).
	\end{aligned}
	\]
	因此
	\[
	\begin{aligned}
	\frac{\partial f}{\partial\vect v}(x_0,y_0)
	&=
	\lim_{t\to0}
	\frac{f(x_0+t\cos\alpha,y_0+t\sin\alpha)-f(x_0,y_0)}{t}\\
	&=
	\frac{\partial f}{\partial x}(x_0,y_0)\cos\alpha
	+
	\frac{\partial f}{\partial y}(x_0,y_0)\sin\alpha.
	\end{aligned}
	\]
\end{proof}

总而言之，对于多元函数的偏导数、方向导数与可微性，有如下关系：
\begin{enumerate}
	\item 可偏导未必连续，可偏导也未必可微；
	\item 沿任意方向的方向导数均存在时，各偏导数必然存在，但函数仍未必可微；
	\item 可微必连续且必可偏导.若 $z=f(\vect x)$，其中
	$\vect x=(x_1,\cdots,x_n)$，则
	\[
	\dd z
	=
	\frac{\partial f}{\partial x_1}\dd x_1
	+\cdots+
	\frac{\partial f}{\partial x_n}\dd x_n
	=
	\frac{\partial f}{\partial\vect x}\cdot\dd\vect x;
	\]
	\item 可微则沿任意方向的方向导数均存在.
\end{enumerate}

\begin{theorem}
	设 $D\subset\R^2$ 为开集，函数 $z=f(x,y)$ 在 $(x_0,y_0)$ 的某邻域内可偏导.若
	\[
	\frac{\partial f}{\partial x},\qquad
	\frac{\partial f}{\partial y}
	\]
	在 $(x_0,y_0)$ 处连续，则 $f$ 在 $(x_0,y_0)$ 处可微.
\end{theorem}

\textbf{注：}若 $\dfrac{\partial f}{\partial x}$、$\dfrac{\partial f}{\partial y}$ 在 $D$ 上连续，则称 $f$ 在 $D$ 上连续可微，记作 $f\in C^1(D)$.上面的定理说明，一阶偏导数连续可以推出可微.

\begin{proof}
	记
	\[
	\Delta z=f(x_0+\Delta x,y_0+\Delta y)-f(x_0,y_0).
	\]
	将其分解为
	\[
	\begin{aligned}
	\Delta z
	={}&f(x_0+\Delta x,y_0+\Delta y)-f(x_0,y_0+\Delta y)\\
	&+f(x_0,y_0+\Delta y)-f(x_0,y_0).
	\end{aligned}
	\]
	分别对两个差值应用一元函数的微分中值定理，存在
	$\theta_1,\theta_2\in(0,1)$，使得
	\[
	\Delta z
	=
	f_x(x_0+\theta_1\Delta x,y_0+\Delta y)\Delta x
	+
	f_y(x_0,y_0+\theta_2\Delta y)\Delta y.
	\]
	由 $f_x,f_y$ 在 $(x_0,y_0)$ 处连续，当
	$\Delta x\to0,\Delta y\to0$ 时，可写成
	\[
	\begin{aligned}
	f_x(x_0+\theta_1\Delta x,y_0+\Delta y)
		&=f_x(x_0,y_0)+\rho(\Delta x,\Delta y),\\
	f_y(x_0,y_0+\theta_2\Delta y)
		&=f_y(x_0,y_0)+\eta(\Delta x,\Delta y),
	\end{aligned}
	\]
	其中 $\rho(\Delta x,\Delta y)\to0$，$\eta(\Delta x,\Delta y)\to0$.代入得
	\[
	\Delta z
	=
	f_x(x_0,y_0)\Delta x+f_y(x_0,y_0)\Delta y
	+\rho\Delta x+\eta\Delta y.
	\]
	又由 Cauchy--Schwarz 不等式，
	\[
	\left|
	\frac{\rho\Delta x+\eta\Delta y}
	{\sqrt{\Delta x^2+\Delta y^2}}
	\right|
	\leq
	\sqrt{\rho^2+\eta^2}\longrightarrow0.
	\]
	故
	\[
	\dd z
	=
	\frac{\partial f}{\partial x}(x_0,y_0)\dd x
	+
	\frac{\partial f}{\partial y}(x_0,y_0)\dd y,
	\]
	即 $f$ 在 $(x_0,y_0)$ 处可微.
\end{proof}

\subsection*{梯度}

设 $z=f(x,y)$ 在 $(x_0,y_0)$ 处可微，则它沿任意方向
$\vect v=(\cos\alpha,\sin\alpha)$ 的方向导数均存在，并且
\[
\begin{aligned}
\frac{\partial f}{\partial\vect v}
&=
\frac{\partial f}{\partial x}\cos\alpha
+
\frac{\partial f}{\partial y}\sin\alpha\\
&=
\left(
\frac{\partial f}{\partial x}\vect i
+
\frac{\partial f}{\partial y}\vect j
\right)
\cdot
\left(\cos\alpha\,\vect i+\sin\alpha\,\vect j\right).
\end{aligned}
\]
由 Cauchy--Schwarz 不等式，
\[
\left|
\frac{\partial f}{\partial\vect v}
\right|
\leq
\left|
\frac{\partial f}{\partial x}\vect i
+
\frac{\partial f}{\partial y}\vect j
\right|
\left|\cos\alpha\,\vect i+\sin\alpha\,\vect j\right|
=
\sqrt{
\left(\frac{\partial f}{\partial x}\right)^2
+
\left(\frac{\partial f}{\partial y}\right)^2
}.
\]
当 $\vect v$ 与
$\dfrac{\partial f}{\partial x}\vect i+
\dfrac{\partial f}{\partial y}\vect j$
同向或反向时，$\left|\dfrac{\partial f}{\partial\vect v}\right|$ 取得最大值.

\begin{definition}
	称向量
	\[
	\operatorname{grad}f
	=
	\frac{\partial f}{\partial x}\vect i
	+
	\frac{\partial f}{\partial y}\vect j
	\]
	为函数 $f$ 的\textbf{梯度}，也记作 $\nabla f$.
\end{definition}

梯度给出了函数变化率最大的方向.具体而言：
\begin{enumerate}
	\item 当 $\operatorname{grad}f\ne\vect 0$ 时，$f$ 沿
	$\operatorname{grad}f$ 的方向增长最快，沿
	$-\operatorname{grad}f$ 的方向减小最快；当
	$\operatorname{grad}f=\vect 0$ 时，各方向的方向导数均为零；
	\item 当 $\operatorname{grad}f\ne\vect 0$ 时，梯度与函数的等值线正交；
	\item $\operatorname{grad}f$ 也可记为 $\nabla f$.
\end{enumerate}

\subsection*{高阶偏导数}

以二元函数 $z=f(x,y)$ 为例.若
$\dfrac{\partial z}{\partial x}$、$\dfrac{\partial z}{\partial y}$
仍然可偏导，则可进一步求二阶偏导数：
\[
\frac{\partial}{\partial x}
\left(\frac{\partial z}{\partial x}\right)
=
\frac{\partial^2z}{\partial x^2}
=f_{xx}(x,y),
\qquad
\frac{\partial}{\partial y}
\left(\frac{\partial z}{\partial y}\right)
=
\frac{\partial^2z}{\partial y^2}
=f_{yy}(x,y),
\]
\[
\frac{\partial}{\partial x}
\left(\frac{\partial z}{\partial y}\right)
=
\frac{\partial^2z}{\partial x\partial y}
=f_{yx}(x,y),
\qquad
\frac{\partial}{\partial y}
\left(\frac{\partial z}{\partial x}\right)
=
\frac{\partial^2z}{\partial y\partial x}
=f_{xy}(x,y).
\]
其中 $f_{xx}$、$f_{yy}$ 称为\textbf{二阶偏导数}，
$f_{xy}$、$f_{yx}$ 称为\textbf{二阶混合偏导数}.

\begin{example}
	混合偏导数的求导次序一般不能任意交换.令
	\[
	f(x,y)=
	\begin{cases}
	xy\dfrac{x^2-y^2}{x^2+y^2},&x^2+y^2\ne0,\\[6pt]
	0,&x^2+y^2=0.
	\end{cases}
	\]
	当 $x^2+y^2\ne0$ 时，
	\[
	\frac{\partial f}{\partial x}
	=
	y\frac{x^4+4x^2y^2-y^4}{(x^2+y^2)^2},
	\qquad
	\frac{\partial f}{\partial y}
	=
	x\frac{x^4-4x^2y^2-y^4}{(x^2+y^2)^2},
	\]
	而在原点处，两者均为零.因此
	\[
	f_{xy}(0,0)
	=
	\lim_{\Delta y\to0}
	\frac{f_x(0,\Delta y)-f_x(0,0)}{\Delta y}
	=-1,
	\]
	\[
	f_{yx}(0,0)
	=
	\lim_{\Delta x\to0}
	\frac{f_y(\Delta x,0)-f_y(0,0)}{\Delta x}
	=1.
	\]
	故 $f_{xy}(0,0)\ne f_{yx}(0,0)$.
\end{example}

\begin{theorem}
	设 $D\subset\R^2$ 为开集，函数 $f:D\to\R$ 的混合偏导数
	$f_{xy}$ 与 $f_{yx}$ 在 $(x_0,y_0)$ 的某邻域内存在，并且二者在
	$(x_0,y_0)$ 处连续，则
	\[
	f_{xy}(x_0,y_0)=f_{yx}(x_0,y_0).
	\]
\end{theorem}

\begin{proof}
	对充分小的非零增量 $\Delta x,\Delta y$，记
	\[
	I(\Delta x,\Delta y)
	=
	\frac{
	f(x_0+\Delta x,y_0+\Delta y)-f(x_0+\Delta x,y_0)
	-f(x_0,y_0+\Delta y)+f(x_0,y_0)}
	{\Delta x\Delta y}.
	\]
	先将分子视为关于 $x$ 的差，再对 $y$ 应用微分中值定理，存在
	$\xi$ 介于 $x_0$ 与 $x_0+\Delta x$ 之间、$\eta$ 介于
	$y_0$ 与 $y_0+\Delta y$ 之间，使得
	\[
	I(\Delta x,\Delta y)=f_{xy}(\xi,\eta).
	\]
	交换两次应用中值定理的顺序，同样存在相应的
	$\xi',\eta'$，使得
	\[
	I(\Delta x,\Delta y)=f_{yx}(\xi',\eta').
	\]
	当 $\Delta x\to0,\Delta y\to0$ 时，
	$(\xi,\eta)\to(x_0,y_0)$ 且
	$(\xi',\eta')\to(x_0,y_0)$.由 $f_{xy}$、$f_{yx}$
	在 $(x_0,y_0)$ 处连续，
	\[
	\lim_{\substack{\Delta x\to0\\\Delta y\to0}}
	I(\Delta x,\Delta y)
	=f_{xy}(x_0,y_0)
	=f_{yx}(x_0,y_0).
	\]
\end{proof}

更一般地，高阶混合偏导数在连续的条件下，其求导次序可以任意交换.例如，若
$f_{xyy}$、$f_{yxy}$、$f_{yyx}$ 均连续，则它们相等，可统一记作
\[
\frac{\partial^3f}{\partial x\partial y^2}.
\]

\subsection*{高阶微分}

设 $D\subset\R^2$ 为开区域，函数 $f:D\to\R$ 连续可微.对于
$z=f(x,y)$，其全微分为
\[
\dd z=\frac{\partial f}{\partial x}\dd x
+\frac{\partial f}{\partial y}\dd y.
\]
若 $\dfrac{\partial f}{\partial x}$、$\dfrac{\partial f}{\partial y}$
也连续可微，则称 $z=f(x,y)$ 二阶连续可微，记为$C^2$,并定义其二阶微分为
\[
\begin{aligned}
\dd^2z
&=\dd\left(
\frac{\partial f}{\partial x}\dd x
+\frac{\partial f}{\partial y}\dd y
\right)\\
&=\dd\left(\frac{\partial f}{\partial x}\right)\dd x
+\frac{\partial f}{\partial x}\dd(\dd x)
+\dd\left(\frac{\partial f}{\partial y}\right)\dd y
+\frac{\partial f}{\partial y}\dd(\dd y)\\
&=
\left(
\frac{\partial^2f}{\partial x^2}\dd x
+\frac{\partial^2f}{\partial y\partial x}\dd y
\right)\dd x
+
\left(
\frac{\partial^2f}{\partial x\partial y}\dd x
+\frac{\partial^2f}{\partial y^2}\dd y
\right)\dd y\\
&=
\left(
\dd x\frac{\partial}{\partial x}
+\dd y\frac{\partial}{\partial y}
\right)^2f.
\end{aligned}
\]

进一步地，若 $f\in C^k(D)$，即 $f$ 的所有不高于 $k$ 阶的偏导数
均连续，可以归纳的证明.
\[
\dd^kz
=
\left(
\dd x\frac{\partial}{\partial x}
+\dd y\frac{\partial}{\partial y}
\right)^kf.
\]

\subsection*{向量值函数}

\begin{definition}
	设
	\[
	\vect F(x_1,\cdots,x_n)
	=
	\begin{pmatrix}
	f_1(x_1,\cdots,x_n)\\
	f_2(x_1,\cdots,x_n)\\
	\vdots\\
	f_m(x_1,\cdots,x_n)
	\end{pmatrix},
	\qquad
	\vect F:D\subset\R^n\to\R^m.
	\]
	称 $\vect F$ 为一个\textbf{向量值函数}，称
	$f_1,\cdots,f_m$ 为 $\vect F$ 的分量函数.
\end{definition}

向量函数 $\vect F$ 连续，当且仅当它的每个分量函数
$f_1,\cdots,f_m$ 都连续.对于 $i=1,\cdots,n$,考虑极限
\[\frac{\partial \vect F}{\partial x_1}=\lim\limits_{\Delta x_1\to 0 }\frac{\vect F(x_1+\Delta x_1,\cdots,x_n)-\vect F(x_1,\cdots,x_n)}{\Delta x_1}\]
若此极限存在且有限，就称$ \vect F$关于$x_1$可偏导.考虑$\vect F $是向量值函数.有

\[
\frac{\partial\vect F}{\partial x_i}
=
\begin{pmatrix}
\dfrac{\partial f_1}{\partial x_i}\\[4pt]
\dfrac{\partial f_2}{\partial x_i}\\
\vdots\\
\dfrac{\partial f_m}{\partial x_i}
\end{pmatrix}.
\]
因此，$\dfrac{\partial\vect F}{\partial x_i}$ 存在，当且仅当所有分量
函数关于 $x_i$ 的偏导数均存在.

\begin{definition}
	向量函数 $\vect F$ 的 \textbf{Jacobi 矩阵}定义为
	\[
	J_{\vect F}
	=
	\frac{D(f_1,\cdots,f_m)}
	{D(x_1,\cdots,x_n)}
	=
	\begin{pmatrix}
	\dfrac{\partial f_1}{\partial x_1}
	&\dfrac{\partial f_1}{\partial x_2}
	&\cdots
	&\dfrac{\partial f_1}{\partial x_n}\\[5pt]
	\dfrac{\partial f_2}{\partial x_1}
	&\dfrac{\partial f_2}{\partial x_2}
	&\cdots
	&\dfrac{\partial f_2}{\partial x_n}\\
	\vdots&\vdots&\ddots&\vdots\\
	\dfrac{\partial f_m}{\partial x_1}
	&\dfrac{\partial f_m}{\partial x_2}
	&\cdots
	&\dfrac{\partial f_m}{\partial x_n}
	\end{pmatrix}_{m\times n}.
	\]
	当 $m=n$ 时，称
	\[
	\det J_{\vect F}
	=
	\frac{\partial(f_1,\cdots,f_n)}
	{\partial(x_1,\cdots,x_n)}
	\]
	为 $f_1,\cdots,f_n$ 关于 $x_1,\cdots,x_n$ 的
	\textbf{Jacobi 行列式}.
\end{definition}

\textbf{注：}
\begin{enumerate}
	\item 交换两个分量函数 $f_i,f_j$，相当于交换 Jacobi 矩阵的
	第 $i,j$ 两行；
	\item 交换两个自变量 $x_i,x_j$，相当于交换 Jacobi 矩阵的
	第 $i,j$ 两列.
\end{enumerate}

\begin{definition}
	设 $\vect x\in D$.若存在一个 $m\times n$ 矩阵 $A$，使得
	\[
	\vect F(\vect x+\Delta\vect x)-\vect F(\vect x)
	=
	A\Delta\vect x+o\bigl(\norm{\Delta\vect x}\bigr),
	\]
	则称 $\vect F$ 在 $\vect x$ 处\textbf{可微}.此时
	\[
	A
	=
	\frac{D(f_1,\cdots,f_m)}
	{D(x_1,\cdots,x_n)}
	=
	J_{\vect F}.
	\]
\end{definition}

与标量函数相同，向量函数可偏导不一定连续，也不一定可微；
可微必连续，并且必可偏导.

